% SOURCES: draft/Features/features.md (Tier-1 capabilities); docs/0.2.0v/plan.md
%          (§11 Feature Checklist, §12 Additional In-Scope, §5 Architecture, §14, §15);
%          README.md; docs/1.0.0v/supervisor-report.md

\chapter{Requirements Gathering}
\label{ch:requirements}

Requirements were elicited from the Tier-1 feature specification
(\texttt{draft/Features/features.md}) and the architecture and constraints in the
execution plan. They are presented here as functional requirements, non-functional
requirements, and the primary use cases; the use-case \emph{diagram} that
accompanies these descriptions appears in Chapter~\ref{ch:analysis},
Figure~\ref{fig:usecase}.

\section{Functional Requirements}
\begin{enumerate}[label=\textbf{FR\arabic*.}, leftmargin=3.4em]
  \item The system shall let a user insert, select, move, reorder, group, and
        delete elements on a canvas, with full undo and redo.
  \item The system shall support the primitive element types --- container, text,
        image, button, link, icon, list, and divider --- and \emph{presets} that
        compose those primitives into larger structures (hero, cards grid,
        navigation, footer, and so on) without introducing new element types.
  \item The system shall let a user define design tokens (colour, spacing, font
        size, font family, line height, radius, shadow) and bind element styles to
        them, with coordinated light and dark theme values.
  \item The system shall support per-breakpoint style values (desktop, tablet,
        mobile, small) and per-state overrides (hover, focus-visible, active) on
        every element.
  \item The system shall create an element from a rectangle the user draws on the
        canvas, proposing the most likely element kind with a confidence hint
        (draw-to-create).
  \item The system shall match a rough drawn layout against bundled professional
        designs and let the user adopt the best match as a full document
        (match-layout).
  \item The system shall validate the document and export a portable HTML/CSS/JS
        bundle as a ZIP archive, together with SEO assets (sitemap, robots).
  \item The system shall expose the document, generation, and export pipeline over
        a headless Model Context Protocol (MCP) server, so the builder can be
        driven without its graphical interface.
  \item The system shall save and open projects as versioned \texttt{.dtw} files,
        with autosave, crash recovery, and reload on external change.
\end{enumerate}

\section{Non-Functional Requirements}
% SOURCES: plan.md §14 Performance Budgets, §15 Security; CLAUDE.md (constraints)
\begin{enumerate}[label=\textbf{NFR\arabic*.}, leftmargin=3.8em]
  \item \textbf{Accessibility.} Export shall be blocked on any critical or serious
        \texttt{axe-core} violation, and the output shall satisfy the project's
        WCAG-aligned rules (single \texttt{<h1>}, no heading skips, \texttt{alt}
        text, ARIA labels on icon-only controls, visible focus, reduced-motion
        support) \cite{wcag21}.
  \item \textbf{Determinism.} Identical input shall yield byte-identical output
        (no timestamps, no random identifiers), so exports are reproducible and
        diff-able.
  \item \textbf{Performance.} Data-layer operations (save, open, undo/redo) and the
        export pipeline shall meet defined budgets (see
        Chapter~\ref{ch:testing}).
  \item \textbf{Security.} The Electron application shall run with context
        isolation, no Node integration, and sandboxing; a strict Content Security
        Policy in production; and Zod-validated IPC boundaries with path
        sanitisation and size caps.
  \item \textbf{Portability.} The exported output shall have no build step and no
        mandatory runtime dependency, so it can be hosted anywhere as static files.
  \item \textbf{Maintainability.} The codebase shall be strict TypeScript with no
        \texttt{any}, named exports, and JSDoc on every exported function in the
        core modules.
\end{enumerate}

\section{Primary Use Cases}
Each use case below is given as a short structured description (actor,
precondition, main flow, postcondition). Together they define the actors and
system boundary drawn in Figure~\ref{fig:usecase}.

\subsection*{UC-1: Author and export a page}
\textbf{Actor:} Author. \quad\textbf{Precondition:} the application is open.
\textbf{Main flow:} the author opens a template or a blank page; inserts and edits
elements on the canvas; adjusts tokens, theme, and per-breakpoint values; runs the
live validation console to resolve any issues; and requests an export.
\textbf{Postcondition:} a valid, accessibility-gated ZIP bundle is written to disk.

\subsection*{UC-2: Draw to create an element}
\textbf{Actor:} Author. \quad\textbf{Precondition:} a container is selected.
\textbf{Main flow:} the author drags a rectangle; the system normalises the box,
ranks the likely element kind with a confidence hint, and offers alternatives; the
author accepts a kind; the element is placed through an ordinary insert operation.
\textbf{Postcondition:} a new element exists in the document at a grid-snapped
position; the drawn pixels never enter the store.

\subsection*{UC-3: Match my layout}
\textbf{Actor:} Author. \quad\textbf{Precondition:} a rough layout has been sketched.
\textbf{Main flow:} the author opens Match; the system computes a structural
fingerprint of the current tree and ranks the bundled professional designs with a
per-dimension breakdown; the author adopts one. \textbf{Postcondition:} the chosen
design is hydrated as an ordinary document through the same path a project load
uses.

\subsection*{UC-4: Drive the builder headlessly (MCP)}
\textbf{Actor:} MCP client / agent. \quad\textbf{Precondition:} the MCP server is
running. \textbf{Main flow:} the client calls \texttt{create\_document}, then
inserts presets, sets tokens, theme, runtime flags, and SEO, then calls
\texttt{run\_a11y\_check} and \texttt{export\_site} --- all without the graphical
interface. \textbf{Postcondition:} an exported bundle that has passed the same
accessibility gate as an interactive export.

\subsection*{UC-5: Save, open, and recover a project}
\textbf{Actor:} Author. \quad\textbf{Precondition:} a project exists in memory or
on disk. \textbf{Main flow:} the author saves the project as a versioned
\texttt{.dtw} file, or opens one; on load, the file passes through Zod validation,
version migration, and validation again; autosave protects unsaved work and crash
recovery restores it after an abnormal exit. \textbf{Postcondition:} the document
in memory matches the persisted state (or the most recent recoverable state).
